\section{Native pilot points, multiplicative families, and transported ideals}
\label{sec:native-pilot-dictionary}

We now distinguish three local inputs: a specified point, the image of
multiplication on a log lattice, and a principal maximal-order ideal
before transport. There are positive comparisons between these inputs,
even for nongeometric Galois arrows. In the attained examples the
point and pre-ideal hulls coincide, while saturation to the whole
multiplication image produces a strictly larger hull.
All the constructions in this section use one native carrier and one
normalization; they do not assume the cross-arithmeticoid comparison
in \cite[Theorem~6.10.1]{JoshiIV}.

\subsection{The local objects and the logarithmic scale}

Let $p>2$ and let $E/\Q_p$ be finite Galois, of degree $d$ and
ramification index $2\le e\le p-2$. Normalize $v(p)=1$, choose a
uniformizer $\pi$, and write
\begin{equation}\label{eq:np-lattice}
 \kappa=\frac{e-1}{e},\qquad
 I=p^{-1}\log\OO_E^\times
   =p^{-1}\mathfrak m_E=\pi^{1-e}\OO_E,\qquad
 u=p^{-1}\log(1+p).
\end{equation}
Logarithm and exponential are inverse isometries on
$1+\mathfrak m_E$ and $\mathfrak m_E$, since
$1/e>1/(p-1)$. The tame different is
$\mathfrak D_{E/\Q_p}=\pi^{e-1}\OO_E$, so $I=\mathfrak D^{-1}$.
The convention $(2p)^{-1}\log\OO_E^\times$ differs only by a
$\Z_p^\times$-factor, which does not affect the ideals below.

For $m\ge1$, set
\[
 T_m=E^{\otimes_{\Q_p}m},\qquad
 A_m=\OO_E^{\otimes_{\Z_p}m},\qquad
 B_m=\text{the normalization of }A_m\text{ in }T_m.
\]
The algebra $T_m$ is a product of $d^{m-1}$ copies of $E$, and
$B_m$ is the product of their integer rings. An ideal written
$\pi^tB_m$ means the product ideal of valuation $t/e$ in every
component, without requiring identical component generators.

Let $\Gamma$ be a family containing the identity of actual integral
Kummer actions on $I$, or more generally any such family in
$\operatorname{Aut}_{\Z_p}(I)$. Each action extends $\Q_p$-linearly
to $E$. We do not suppose that the actual Galois family exhausts
this general linear group. With the same arrow in each factor, put
\begin{equation}\label{eq:np-hull}
 \mathcal H_\Gamma(S)=
 \overline{\Span_{B_m}
 \{Fx_1\otimes\cdots\otimes Fx_m:
             F\in\Gamma,\ (x_1,\ldots,x_m)\in S\}}
 \qquad(S\subset E^m).
\end{equation}
This is an orbit followed by a module hull. Taking a product convex
hull before tensoring is a separate operation.

Write $\operatorname{Kum}$ for the integral Kummer map on principal
units. The standard Bloch--Kato exponential, extended by
$p^{-n}\operatorname{Kum}(\exp(p^nx))$ when necessary, has inverse
$\log_{\rm BK}^{\rm std}$ satisfying
\[
 \log_{\rm BK}^{\rm std}(\operatorname{Kum}(1+pa))
       =\log(1+pa).
\]
The following coordinate map is different:
\begin{equation}\label{eq:np-rho}
 \rho=p^{-1}\log_{\rm BK}^{\rm std},\qquad
 \lambda_p(a)=\rho(\operatorname{Kum}(1+pa))
             =p^{-1}\log(1+pa).
\end{equation}
The three displays \cite[(9.7.1.1), (9.7.2.1)--(9.7.2.2),
PDF p.~110]{JoshiIII} require this distinction. If the exponential
has its standard cohomological scalar meaning and the class is the
unscaled class of $1+pa$ defined on p.~109, its inverse cannot take
that class to $\lambda_p(a)$. Indeed, for nonzero integral $a$,
\[
 \exp_{\rm BK}(\lambda_p(a))
       =p^{-1}\operatorname{Kum}(1+pa)
       \ne\operatorname{Kum}(1+pa).
\]
The Kummer class is nonzero because its logarithm has nonzero
leading term $pa$. Formula~\eqref{eq:np-rho} is a consistent
normalized coordinate, but it is not that inverse formula.

Passing from $\rho$ to $\log_{\rm BK}^{\rm std}$ multiplies each
entry, including each background entry, by $p$. Since integral
Kummer arrows commute with this scalar, every length-$m$ hull satisfies
\begin{equation}\label{eq:np-standard-scale}
 H^{\rm std}=p^mH^\rho,\qquad
 V_B(H^{\rm std})=V_B(H^\rho)-m\log p,
 \quad
 V_B(H)=d^{-m}\log\mu_B(H),\quad \mu_B(B_m)=1.
\end{equation}
The determinant of multiplication by $p^m$ on $T_m$ is $p^{md^m}$,
which proves the volume formula. If instead $\mu_A(A_m)=1$, then
$V_A(H)=V_B(H)+d^{-m}\log[B_m:A_m]$. This constant cancels when
both sets use the same reference order.

\subsection{A point bridge that survives every integral arrow}

\begin{proposition}[Transported root and crystalline point]
\label{prop:np-point-bridge}
For $0\ne a\in\OO_E$, set $r=v(a)$ and
$k=\lfloor r+\kappa\rfloor$. Every
$F\in\operatorname{Aut}_{\Z_p}(I)$ satisfies
\begin{equation}\label{eq:np-point-val}
             v(F\lambda_p(a))=v(Fa),\qquad v(Fu)=v(F1).
\end{equation}
Consequently
\begin{equation}\label{eq:np-point-hulls}
 \mathcal H_\Gamma(\{(\lambda_p(a),u,\ldots,u)\})
      =\mathcal H_\Gamma(\{(a,1,\ldots,1)\}).
\end{equation}
\end{proposition}
\begin{proof}
The choice of $k$ gives
$a\in p^kI\setminus p^{k+1}I$. In the convergent expansion
\[
 \lambda_p(a)-a
       =\sum_{n\ge2}(-1)^{n+1}\frac{p^{n-1}a^n}{n},
\]
every term has valuation at least $1+2r$: for $p>2$ and $n\ge2$,
$v_p(n)\le n-2$. Hence
\[
 v(\lambda_p(a)-a)\ge1+2r\ge k+1-\kappa,\qquad
              \lambda_p(a)-a\in p^{k+1}I.
\]
As $F$ preserves each $p^nI$ in both directions,
\[
 v(Fa)<k+1-\kappa,\qquad
 v(F(\lambda_p(a)-a))\ge k+1-\kappa.
\]
The strict ultrametric inequality proves the first equality in
\eqref{eq:np-point-val}; apply it to $a=1$ for the second.

In every component of $T_m$, the two nonzero tensors
$F\lambda_p(a)\otimes(Fu)^{\otimes(m-1)}$ and
$Fa\otimes(F1)^{\otimes(m-1)}$ have the same valuation.
Their componentwise ratio is a unit of $B_m$, so they generate
the same principal $B_m$-ideal. Taking spans and closures proves
\eqref{eq:np-point-hulls}. No $E$-linearity or multiplicativity
of $F$ was used.
\end{proof}

The same proof allows different prescribed integral arrows in
different factors. Root-of-unity ambiguities can be handled by
applying the proposition separately to each permitted $\zeta a$
and then taking their union; one does not assert
$F(\zeta a)=\zeta F(a)$. Also, \eqref{eq:np-point-hulls} is a
statement about module hulls of raw orbits. The tensor map is not
additive in its tuple of inputs, so equality after a product
convex hull needs a further argument. An attaining-point argument
below supplies it in the explicit example.

\subsection{A whole multiplication image requires a whole unit family}

\begin{proposition}[Exact saturated Kummer family]
\label{prop:np-saturated}
For $0\ne a\in\OO_E$, the group
$U_a=1+paI=1+a\mathfrak m_E$ satisfies, as actual sets,
\begin{equation}\label{eq:np-saturated}
 \rho(\operatorname{Kum}(U_a))=aI,\qquad
 \log_{\rm BK}^{\rm std}(\operatorname{Kum}(U_a))=paI.
\end{equation}
\end{proposition}
\begin{proof}
The ideal $paI=a\mathfrak m_E$ lies in the common strict convergence
disk of logarithm and exponential. These maps preserve valuations,
so $\log(1+paI)=paI$. More explicitly, for every $z\in aI$,
$\exp(pz)\in U_a$ and
$\rho(\operatorname{Kum}(\exp(pz)))=z$. The inverse series gives
the reverse containment. The positive-valuation ideal also makes
$U_a$ closed under multiplication and inversion.
\end{proof}

In particular, the background family is the whole principal-unit
group $U_1$, not the singleton $1+p$. Taking products in
\eqref{eq:np-saturated} realizes $aI\times I^{m-1}$ before
any arrows or hulls are applied. This equality therefore survives
the same transport, product convex closure, tensor map, and final
module hull, in any specified order.

\begin{proposition}[The singleton can be strictly smaller]
\label{prop:np-strict}
Put
\[
 P(a)=\mathcal H_\Gamma(\{(\lambda_p(a),u,\ldots,u)\}),\qquad
 S(a)=\mathcal H_\Gamma(aI\times I^{m-1}).
\]
If $v(a)<\lfloor v(a)+\kappa\rfloor$, then $P(a)\subsetneq S(a)$.
The strict inclusion persists if both product inputs first undergo
closed $\Z_p$-convex closure.
\end{proposition}
\begin{proof}
Writing $r=v(a)$ and $k=\lfloor r+\kappa\rfloor$, the logarithmic
expansion gives $\lambda_p(a)/a\in\OO_E^\times$. Also
$u\in\OO_E\subset I$, so the singleton input lies in the whole
input. Every transported singleton lies in the closed product
lattice $p^kI\times I^{m-1}$, so every tensor in $P(a)$ has
component valuation at least $k-m\kappa$.
The whole input at the identity contains
\[
                   (a\pi^{1-e},\pi^{1-e},\ldots,\pi^{1-e}).
\]
Its tensor has component valuation $r-m\kappa<k-m\kappa$,
proving strictness. Convex closure keeps the smaller input in its
closed product container and retains the displayed larger input.
\end{proof}

Proposition~\ref{prop:np-saturated} supplies a positive cohomological
realization of a specified multiplication image. It is an enlargement
of a single-class recipe. Proposition~\ref{prop:np-strict} shows
why equality of coefficient norms alone cannot justify omitting
this enlargement.

\subsection{One actual arrow for six normalized inputs}

We use the field of \eqref{eq:gl-139-field}: $p=139$,
$v(b_0)=1$, $\pi^{105}=b_0$, and
$E=\Q_p(\mu_{210},\pi)$. Thus $(d,e,f)=(210,105,2)$ and
$I=\pi^{-104}\OO_E$. Put
\[
 r_0=\pi^{15},\quad \alpha_s=r_0^s,\quad
 \tau_s=\lambda_p(\alpha_s),\quad
 n_s=\lfloor s/7\rfloor,\quad
 k_s=\lfloor s/7+104/105\rfloor
 \qquad(s=1,4,9),
\]
and define two normalized families
\begin{equation}\label{eq:np-six-vectors}
             T_s=\tau_s/p^{k_s},\qquad
             X_s=\alpha_s\pi^{-104}/p^{n_s}.
\end{equation}

\begin{lemma}[A common arrow for both families]
\label{lem:np-six}
One actual integral Kummer arrow $F$ satisfies
\[
           v(Fu)=v(FT_s)=v(FX_s)=-104/105
                  \qquad(s=1,4,9).
\]
\end{lemma}
\begin{proof}
The triples $(n_s)$ and $(k_s)$ are $(0,0,1)$ and $(1,1,2)$.
The native valuations are
\[
 \begin{array}{c|ccc}
 s&1&4&9\\ \hline
 v(T_s)&-90/105&-45/105&-75/105\\
 v(X_s)&-89/105&-44/105&-74/105 .
 \end{array}
\]
Thus all six vectors are in $\pi I\setminus pI$.
The Kummer conjugates show $\operatorname{Tr}(X_s)=0$, since
$105\nmid15s-104$, including for the negative powers occurring
here. The direct norm computation used in
Theorem~\ref{thm:gl-three-labels} gives
\[
 \operatorname{Tr}(T_s)
     =\frac{30\log(1+p^7b_0^s)}{p^{k_s+1}},\qquad
 v(\operatorname{Tr}(T_s))=s+6-k_s\in\{6,9,13\},
 \qquad\operatorname{Tr}(u)=210u\in\Z_p^\times.
\]
For the norm formula, the seven distinct conjugates of
$1+p\pi^{15s}$ multiply to $1+p^7b_0^s$, and there are
fifteen repetitions over the unramified quadratic subfield.

Let $V=I/pI$ and $L=\pi I/pI$, of codimension two.
Lemma~\ref{lem:gl-integral-basis} gives coordinates
$V=\mathbf F_p a\oplus\mathbf F_p b\oplus W$ in which
the $b$-coordinate is a unit multiple of the trace.
The trace is onto on $\pi I$: the element $1/210$ lies in
$\OO_E\subset\pi I$ and has trace one. Consequently
$\ker B\to V/L$ is onto. The vector $u$ is in $L$ with nonzero
$b$-coordinate, and the six other vectors are nonzero in
$L\cap\ker B$.

There is an actual field automorphism
$\sigma(\pi)=\zeta_7\pi$, fixing $\Q_p(\mu_{210})$.
Modulo $pI$, the leading term of $T_s$ is
$\alpha_s/p^{k_s}$. Proposition~\ref{prop:np-point-bridge}
places its logarithmic error in $pI$. The characters on
$T_s$ and $X_s$ are therefore respectively
$\zeta_7^s$ and $\zeta_7^{15s-104}=\zeta_7^{s+1}$.
All six are nontrivial. Since
$\mu_7\cap\mathbf F_{139}^\times=\{1\}$, each of their
projective orbits has length seven. Indeed a difference
between a nontrivial such scalar and a scalar of
$\mathbf F_{139}^\times$ is a residue-field unit and
cannot annihilate a nonzero vector.

For each vector there is at most one
$h\in\{0,\ldots,6\}$ placing $\sigma^h$ of that vector on
the line $\mathbf F_p a$. Avoid all six forbidden choices
at once. This common inertia action fixes $u$ and preserves
$L$ and the trace. Now the six vectors satisfy all hypotheses
of Lemma~\ref{lem:gl-finite-family}, with $r=6$ and
$r+1<139$. Its full-Galois lift sends them and $u$ outside
$L$ simultaneously. Finally
Proposition~\ref{prop:gl-variance} converts the canonical
action $M$ into a Kummer action by taking the inverse
Galois arrow: the resulting action is a $\Z_p^\times$-multiple
of $M$. This preserves all seven minimum valuations.
\end{proof}

\begin{theorem}[Exact point and whole-family hulls]
\label{thm:np-two-hulls}
For $s=j^2$, $j=1,2,3$, and $m=j+1$, define
\[
 P_j^\rho=\mathcal H_\Gamma(\{(\tau_s,u,\ldots,u)\}),\qquad
 S_j^\rho=\mathcal H_\Gamma(\alpha_sI\times I^{m-1}),
\]
where $\Gamma$ is the full integral Galois--Kummer family.
Then
\begin{equation}\label{eq:np-two-hulls}
 P_j^\rho=\pi^{105k_s-104m}B_m,\qquad
 S_j^\rho=\pi^{105n_s-104m}B_m=p^{-1}P_j^\rho .
\end{equation}
One common arrow attains both bounds in all three blocks.
The same equalities hold with intermediate closed product
convex hulls.
\end{theorem}
\begin{proof}
The containments $\tau_s\in p^{k_s}I$ and
$\alpha_sI\subseteq p^{n_s}I$ give componentwise lower
valuation bounds $k_s-m\kappa$ and $n_s-m\kappa$ for
the two transported families. Lemma~\ref{lem:np-six}
attains the first bound on
$(\tau_s,u,\ldots,u)$, and the second on
$(\alpha_s\pi^{-104},u,\ldots,u)$, using the same arrow:
their first entries are respectively $p^{k_s}T_s$ and
$p^{n_s}X_s$. Each attaining tensor has that valuation
in every field component. Its principal $B_m$-span is
the entire asserted product ideal, proving the reverse
containments and closedness. Since $k_s=n_s+1$,
the two ideals differ by $p^{-1}$.

The containing product lattices
$p^{k_s}I\times I^{m-1}$ and $p^{n_s}I\times I^{m-1}$
are closed and convex. Convex closure cannot violate
the upper containments and retains the attaining tuples.
\end{proof}

Proposition~\ref{prop:np-point-bridge} gives the same point
hull with $(\alpha_s,1,\ldots,1)$; the same argument now
also treats its product convex hull. The formulas remain
valid for independent integral factor arrows if the
diagonal witnesses are retained. For reference, the
exponents of $\pi$ in the two logarithmic scales are
\[
 \begin{array}{c|c|rr|rr}
 j&m&P_j^\rho&S_j^\rho&P_j^{\rm std}&S_j^{\rm std}\\ \hline
 1&2&-103&-208&107&2\\
 2&3&-207&-312&108&3\\
 3&4&-206&-311&214&109
 \end{array}
\]
The last two columns add $105m$ to the corresponding
$\rho$ columns. Nonintegrality in a $\rho$ column is
therefore not a violation of a bound stated in the
standard cohomological coordinate.

\subsection{The maximal-order ideal before transport}

For $z=a\otimes1\otimes\cdots\otimes1$ and
$\Phi_F=F^{\otimes m}$, put
\begin{equation}\label{eq:np-preideal}
 \mathcal M_\Gamma(a)=
 \overline{\Span_{B_m}
       \{\Phi_F(b):b\in zB_m,\ F\in\Gamma\}}.
\end{equation}
The ideal $zB_m$ is the input, not an ideal formed after
applying $\Phi_F$.

\begin{proposition}[A trace-dual comparison for one source set]
\label{prop:np-sandwich}
For $0\ne a\in\OO_E$,
\begin{equation}\label{eq:np-sandwich}
 \mathcal H_\Gamma(\{(a,1,\ldots,1)\})
       \subseteq\mathcal M_\Gamma(a)
       \subseteq\mathcal H_\Gamma(aI\times I^{m-1}).
\end{equation}
\end{proposition}
\begin{proof}
The first containment follows from $z\in zB_m$.
Take the integral trace dual of $A_m$ in the finite etale
algebra $T_m$. Trace factorizes on pure tensors; choosing
an integral basis and its trace-dual basis in each factor
therefore gives
\[
 A_m^\vee
    =(\mathfrak D_{E/\Q_p}^{-1})^{\otimes_{\Z_p}m}
    =I^{\otimes_{\Z_p}m}.
\]
If $b\in B_m$ and $c\in A_m\subset B_m$, the product
$bc$ is integral in every component field. Its trace
belongs to $\Z_p$, proving $B_m\subseteq A_m^\vee$.
Consequently $zB_m$ is contained in the $\Z_p$-span
of pure tensors with entries in $aI\times I^{m-1}$.
Apply the $\Q_p$-linear map $\Phi_F$, then take the
common $B_m$-span and closure.
\end{proof}

In particular, the preceding native example gives bounds
for the same pre-hull input:
\begin{equation}\label{eq:np-native-interval}
 P_j^\rho\subseteq\mathcal M_\Gamma(\alpha_{j^2})
       \subseteq S_j^\rho=p^{-1}P_j^\rho,
\end{equation}
and hence
\[
 (m\kappa-k_{j^2})\log p
 \le V_B(\mathcal M_\Gamma(\alpha_{j^2}))
 \le(m\kappa-n_{j^2})\log p.
\]
This is a valid interval bound of width $\log p$.
The next result sharpens it and determines the middle
ideal exactly in the attained examples.

\begin{proposition}[The sharper trace-dual bound]
\label{prop:np-sharp-preideal}
Let $0\ne a\in\OO_E$, $r=v(a)$, and
$k=\lfloor r+\kappa\rfloor$. If a $\Q_p$-linear map
$\Phi:T_m\to T_m$ satisfies
$\Phi(A_m^\vee)\subseteq A_m^\vee$, then
\begin{equation}\label{eq:np-sharp-preideal}
 \Span_{B_m}\Phi(zB_m)
       \subseteq\pi^{ek-m(e-1)}B_m,
 \qquad z=a\otimes1\otimes\cdots\otimes1.
\end{equation}
In particular this bounds $\mathcal M_\Gamma(a)$ for
every family of integral factor arrows used above.
The upper bound itself does not require $\Phi$ to
factor as a tensor product.
\end{proposition}

\begin{proof}
Take all trace duals with respect to the absolute algebra
trace $\operatorname{Tr}_{T_m/\Q_p}$. Since $E/\Q_p$
is Galois, every component field of $T_m$ is a copy of
$E$. The idempotents of $B_m$ isolate the components
of the trace sum, so
\begin{equation}\label{eq:np-two-duals}
 B_m^\vee=\prod_\alpha I_\alpha
       \subseteq A_m^\vee=I^{\otimes_{\Z_p}m},
\end{equation}
where $I_\alpha$ is the inverse different in the
corresponding copy of $E$. There is no additional
degree or component-count factor in this formula.
The containment follows from $A_m\subseteq B_m$
by contravariance of the trace dual.

Each component of $z$ has valuation $r$. Because
$r-k\ge-\kappa$, one has
\[
 p^{-k}zB_m\subseteq B_m^\vee\subseteq A_m^\vee.
\]
It follows that $\Phi(zB_m)\subseteq p^kA_m^\vee$.
Every pure tensor in $I^{\otimes m}$ has component
valuation at least $-m\kappa$, and
$(\pi^{1-e})^{\otimes m}$ attains this value in every
component. Consequently
\[
 \Span_{B_m}A_m^\vee=\pi^{-m(e-1)}B_m.
\]
Multiplying by $p^k$ proves
\eqref{eq:np-sharp-preideal}. The bounding ideal is
closed, so unions and closures preserve this upper
bound. Integral factor arrows preserve
$I^{\otimes m}=A_m^\vee$, as required.
\end{proof}

\begin{corollary}[Exact pre-ideal hull]
\label{cor:np-exact-preideal}
If the point orbit attains depth $k-m\kappa$ in every
component, then
\[
 \mathcal M_\Gamma(a)=
 \mathcal H_\Gamma(\{(a,1,\ldots,1)\})
       =\pi^{ek-m(e-1)}B_m.
\]
In the degree-210 three-label example, this gives
\begin{equation}\label{eq:np-exact-preideal}
 \begin{gathered}
 \mathcal M_\Gamma(\alpha_{j^2})=P_j^\rho,
 \qquad S_j^\rho=p^{-1}\mathcal M_\Gamma(\alpha_{j^2}),\\
 V_B(\mathcal M_\Gamma(\alpha_{j^2}))
       =(m\kappa-k_{j^2})\log p,
 \qquad m=j+1,\quad j=1,2,3.
 \end{gathered}
\end{equation}
The same actual full-Galois arrow can witness all three
point attainments.
\end{corollary}

\begin{proof}
The point-orbit containment in
Proposition~\ref{prop:np-sandwich} and the assumed
attainment give the reverse inclusion to
\eqref{eq:np-sharp-preideal}. In the stated example,
Proposition~\ref{prop:np-point-bridge} identifies each
transported root-point principal ideal with its
normalized crystalline-point ideal, including the
background factors. Theorem~\ref{thm:np-two-hulls}
supplies the common attaining arrow and the separate
whole-family formula. The measure follows from the
depth of the resulting product ideal.
\end{proof}

This comparison has the input type of the containing
calculation in \cite[Proposition~1.2, PDF pp.~10--11;
Theorem~1.10, Step~(v), pp.~27--28]{IUT4}.
Indeed $\Phi_F$ preserves
$I^{\otimes m}$, and hence
$(\log\OO_E^\times)^{\otimes m}=p^mI^{\otimes m}$,
as required there. The displayed input is
$\phi(p^\lambda B_m)$, with
$\lambda=v(\underline q^{\,j^2})$ at a bad place.
Here $\underline q=q^{1/(2\ell)}$ is the root fixed in
\cite[Example~3.2(iv), PDF p.~71]{IUT1}.
In the example $e=105,p=139$, its parameters are
\[
 d_i=104/105,\qquad a_i=1/105,\qquad b_i=-1/105.
\]
Thus its displayed containing ideal has depth
\[
 \lfloor\lambda-d_I-a_I\rfloor-b_I
       =\lfloor\lambda\rfloor-m\kappa,
       \qquad \lambda=j^2/7,
\]
which is the upper endpoint in
\eqref{eq:np-native-interval}. By
\eqref{eq:np-exact-preideal} it is a larger containing
ideal, equal to $p^{-1}$ times the exact pre-ideal hull
in these examples. A loose upper bound is not a
contradiction, and its agreement with the whole-family
formula does not identify those input sets.

If all tensor coordinates are changed from $\rho$
to the standard logarithm, the pre-ideal input changes
from $zB_m$ to $p^mzB_m$. Its hull then becomes
$p^m\mathcal M_\Gamma(\alpha_{j^2})=P_j^{\rm std}$,
with measure shift $-m\log p$. This conversion does
not insert $p^m$ into a printed $\phi(p^\lambda B_m)$
input without changing its coordinate description.

\subsection{The source distinctions that remain}

Three distinctions are necessary when using these local
interfaces in the published constructions.

First, \cite[\S9.6.2 and Theorem-Definition~9.8.1.1,
PDF pp.~109, 114--115]{JoshiIII} defines a chosen class
from $1+p\,q^{1/(2\ell)}$ and background classes from $1+p$.
The optional family in its Remark~9.6.2.1 is printed
with exponent $j/(2\ell)$, whereas the squared-exponent
pilot in \cite[Theorem~3.11; Remark~3.11.1(i)]{IUT3}
involves $\underline q^{\,j^2}$.
The valuation scaling of \cite[(4.2.2.2)]{JoshiIII}
is additional marking data. It does not by itself
replace an algebraic root by its squared power.
For example, in our fixed native field
$\lambda_p(r_0)$ belongs to $pI\setminus p^2I$,
while $\lambda_p(r_0^9)$ belongs to $p^2I\setminus p^3I$.
No integral log-lattice isomorphism identifies these
two specified elements. A change $v\mapsto cv$ changes
$v(p)$ as well and does not change this integral content.

Second, the monoid representation in
\cite[Theorem~3.11(i)(b), PDF pp.~153--154]{IUT3}
is a subset equipped with a multiplicative action;
the pilot also has the fractional-ideal realization of
its Proposition~3.7(v) and Definition~3.8(i).
A point, an operator image, and a principal ideal before
transport have different quantifiers.
Proposition~\ref{prop:np-point-bridge} supplies an equality
for specified point hulls, and
Proposition~\ref{prop:np-saturated} supplies a set equality
for a specified whole family. Neither permits a
$B_m$-span to be moved across a non-$B_m$-linear arrow.
Proposition~\ref{prop:np-sharp-preideal} instead
bounds the entire pre-ideal before the arrow is
applied; together with point attainment it proves
\eqref{eq:np-exact-preideal} without moving a module
span across that arrow.

Third, the standard versus normalized logarithmic map
must be the same on both sides of a volume comparison.
Equation~\eqref{eq:np-standard-scale} is its explicit
conversion, including background factors. The
native fixed-source covariance of an all-isomorphism
collation cannot at the same time be interpreted as
a change of its source element or source principal ideal.
Remark~6.10.2 of \cite{JoshiIV} itself restricts direct
comparison of local quantities at different arithmeticoids.

The degree-210 calculation applies locally to a split
rational Frey curve with full rational $2$-torsion and
$v(q)=2$. Tate uniformization
\cite[Theorem~C.14.1]{Silverman} makes a chosen square root
of $q$ rational: its class is a rational point of order
two, so its Galois sign must be a power of $q$; valuations
force that power to be zero. Thus it supplies precisely
the $b_0$ allowed above, without requiring $q=p^2$.
This local observation does not establish all the
initial data in \cite[Definition~3.1]{IUT1}.

There is also a small-parameter limitation to using the
particular existence proof of \cite[\S5.8]{JoshiIV}.
Its prime threshold must satisfy $\xi_{\rm prm}>10$,
because $\theta(10)=\log210<20/3$. With the normalized
$Q$ of its Definition~5.4.1 and Theorem~5.7.1,
every candidate with $\sqrt Q\le\ell=7$ is therefore
included in the proof's small-$Q$ exceptional set, if
it is in the chosen bounding domain. The opening display
of Section~5.8 omits the degree divisor present in that
definition. If this display is instead read literally
as a different unnormalized quantity, the normalized
value and numerical prime window cannot be reused for it.
Neither reading turns a local example into an automatic
outside-exception global input. This statement concerns
that proof-constructed set, not every possible sufficient
exceptional set or a direct verification of initial data.

Finally, the actual global family, its product weights,
the Ind3 operation, and the cross-Frobenius comparison
are not computed in this section. The trace-dual bound,
the attained exact pre-ideal equality, and the two positive
Kummer interfaces are local theorems
about their explicitly specified sets. Their $p$-adic
analytic and full local-Galois hypotheses have not yet
been completely formalized in Lean. They neither prove
nor disprove the unconditional abc conjecture.
